% Part: first-order-logic
% Chapter: axiomatic-deduction
% Section: proving-things

\documentclass[../../../include/open-logic-section]{subfiles}

\begin{document}

\iftag{FOL}
      {\olfileid{fol}{axd}{pro}}
      {\olfileid{pl}{axd}{pro}}

\olsection{\usetoken{P}{derivation}-এর উদাহরণ}


\begin{ex}
ধরা যাক, আমরা $(\lnot !D \lor !E) \lif (!D \lif !E)$ প্রমাণ করতে
চাই। স্পষ্টতই এটি আমাদের কোনো স্বতঃসিদ্ধের রূপ নয়, তাই এটিকে
!!{derive} করতে \MP{} বিধি ব্যবহার করতে হবে। আমাদের একমাত্র বিধি
মোডাস পোনেন্স; $!A$ ও~$!A \lif !B$ দেওয়া থাকলে এটি~$!B$-কে
সমর্থন করার অনুমতি দেয়। একটি কৌশল হলো \olref[prp]{ax:lor3}-এ
$!A$-এর স্থানে $\lnot !D$, $!B$-র স্থানে $!E$ এবং $!C$-র স্থানে
$!D \lif !E$ নেওয়া, অর্থাৎ এই রূপটি নেওয়া:
\[
(\lnot !D \lif (!D \lif !E)) \lif ((!E \lif (!D \lif !E)) \lif ((\lnot
!D \lor !E) \lif (!D \lif !E))).
\]
কেন? মোডাস পোনেন্স দুবার প্রয়োগ করলে শেষাংশটি পাওয়া যায়, আর সেটিই
আমাদের লক্ষ্য। আবার সহজেই দেখা যায়, $\lnot !D \lif (!D \lif !E)$
হলো \olref[prp]{ax:lnot2}-এর একটি রূপ এবং $!E \lif (!D \lif !E)$
হলো \olref[prp]{ax:lif1}-এর একটি রূপ। তাই আমাদের !!{derivation}:
\begin{derivation}
  1. &  $\lnot !D \lif (!D \lif !E)$ & \olref[prp]{ax:lnot2} \\
  2. & $(\lnot !D \lif (!D \lif !E)) \lif {}$\\
  &\qquad $((!E \lif (!D \lif !E)) \lif ((\lnot !D \lor !E) \lif (!D \lif !E)))$ & \olref[prp]{ax:lor3}\\
  3. & $(!E \lif (!D \lif !E)) \lif ((\lnot !D \lor !E) \lif (!D \lif !E))$ & 1, 2, \MP\\
  4. & $!E \lif (!D \lif !E)$ & \olref[prp]{ax:lif1}\\
  5. & $(\lnot !D \lor !E) \lif (!D \lif !E)$ & 3, 4, \MP
\end{derivation}
\end{ex}

\begin{ex}\ollabel{ex:identity}
$!D \lif !D$-এর কোনো !!a{derivation} খোঁজার চেষ্টা করি। এটি কোনো
স্বতঃসিদ্ধের রূপ নয়, তাই এটিকে !!{derive} করতে \MP{} ব্যবহার করতে
হবে। \olref[prp]{ax:lif1} হলো~$!A \lif !B$ রূপের একটি স্বতঃসিদ্ধ,
যার ওপর~\MP প্রয়োগ করা যায়। অবশ্যই, এই ক্ষেত্রে \MP{} যে~$!B$-কে
সঠিক ধাপ হিসেবে সমর্থন করবে, কাজে লাগতে হলে সেটিকে~$!D \lif !D$
হতে হবে, কারণ এটিই আমরা !!{derive} করতে চাই। তার মানে $!A$-কেও
$!D$ হতে হবে; অর্থাৎ \olref[prp]{ax:lif1}-এর এই রূপটি বিবেচনা করা যায়:
\[
!D \lif (!D \lif !D)
\]
\MP প্রয়োগ করতে হলে সংশ্লিষ্ট দ্বিতীয় পূর্বধারণা, অর্থাৎ~$!A$-কেও
সমর্থন করতে হবে। কিন্তু আমাদের ক্ষেত্রে সেটি~$!D$, আর শুধু~$!D$-কে
আমরা !!{derive} করতে পারব না। তাই ভিন্ন কৌশল দরকার।

শুধু~$\lif$-সংবলিত অন্য স্বতঃসিদ্ধটি হলো \olref[prp]{ax:lif2}, অর্থাৎ
\[
(!A \lif (!B \lif !C)) \lif ((!A \lif !B) \lif (!A \lif !C))
\]
\MP{} দুবার প্রয়োগ করে শেষের বাসাবাঁধা শর্তাধীন সূত্রটিতে পৌঁছানো
যায়। আবার তার মানে, \olref[prp]{ax:lif2}-এর এমন একটি রূপ চাই যেখানে
$!A \lif !C$ হলো $!D \lif !D$, অর্থাৎ আমাদের লক্ষ্য !!{formula}।
তাহলে অবশ্যই $!A$ ও~$!C$ উভয়ই~$!D$। কীভাবে~$!B$ বেছে নিলে $!A
\lif (!B \lif !C)$ ও~$!A \lif !B$, অর্থাৎ আমাদের ক্ষেত্রে $!D
\lif (!B \lif !D)$ ও~$!D \lif !B$-ও !!{derivable} হবে? প্রথমটি
আমরা~$!B$ যাই নিই না কেন, ইতিমধ্যেই \olref[prp]{ax:lif1}-এর একটি
রূপ। আর $!B$ যদি $(!D \lif !D)$ হয়, তবে $!D \lif !B$-ও
\olref[prp]{ax:lif1}-এর আরেকটি রূপ। তাই আমাদের !!{derivation}:
\begin{derivation}
  1. & $!D \lif ((!D \lif !D) \lif !D)$ & \olref[prp]{ax:lif1}\\
  2. & $(!D \lif ((!D \lif !D) \lif !D)) \lif {}$\\
  & \qquad $((!D \lif (!D \lif !D)) \lif (!D \lif !D))$ & \olref[prp]{ax:lif2}\\
  3. & $(!D \lif (!D \lif !D)) \lif (!D \lif !D)$ & 1, 2, \MP\\
  4. & $!D \lif (!D \lif !D)$ & \olref[prp]{ax:lif1}\\
  5. & $!D \lif !D$ & 3, 4, \MP
\end{derivation}
\end{ex}

\begin{ex}\ollabel{ex:chain}
কখনো অন্য কয়েকটি !!{formula}s~$\Gamma$ থেকে কোনো !!{formula}-র
!!a{derivation} আছে দেখাতে চাই। যেমন, দেখাই যে $\Gamma = \{!A \lif
!B, !B \lif !C\}$ থেকে $!A \lif !C$ !!{derive} করা যায়।
\begin{derivation}
  1. & $!A \lif !B$ & \Hyp\\
  2. & $!B \lif !C$ & \Hyp\\
  3. & $(!B \lif !C) \lif (!A \lif (!B \lif !C))$ & \olref[prp]{ax:lif1} \\
  4. & $!A \lif (!B \lif !C)$ & 2, 3, \MP\\
  5. & $(!A \lif (!B \lif !C)) \lif {}$\\
  & \qquad $((!A \lif !B) \lif (!A \lif !C))$ & \olref[prp]{ax:lif2}\\
  6. & $((!A \lif !B) \lif (!A \lif !C))$ & 4, 5, \MP\\
  7. & $!A \lif !C$ & 1, 6, \MP
\end{derivation}
``\Hyp'' (``hypothesis'' অর্থাৎ অনুমান) চিহ্নিত পংক্তিগুলি বোঝায় যে
সেই পংক্তির !!{formula}-টি~$\Gamma$-র একটি !!a{element}।
\end{ex}

\begin{prop}
  \ollabel{prop:chain} $\Gamma \Proves !A \lif !B$ এবং $\Gamma
  \Proves !B \lif !C$ হলে $\Gamma \Proves !A \lif !C$।
\end{prop}

\begin{proof}
  ধরা যাক, $\Gamma \Proves !A \lif !B$ এবং $\Gamma \Proves !B
  \lif !C$। তাহলে $\Gamma$ থেকে $!A \lif !B$-এর কোনো
  !!a{derivation} এবং $\Gamma$ থেকে~$!B \lif !C$-এরও কোনো
  !!a{derivation} আছে। দুটিকে পরপর বসিয়ে একটি !!{derivation} বানাও।
  এবার আগের উদাহরণের !!{derivation}-এর 3--7 নম্বর পংক্তি যোগ করো।
  এটি~$\Gamma$ থেকে $!A \lif !C$—নতুন !!{derivation}-টির শেষ
  পংক্তি—এর কোনো !!a{derivation}। লক্ষ করো, 4 ও~7 নম্বর পংক্তির
  সমর্থন তখনও বৈধ থাকে, যদি 1 নম্বর পংক্তির উল্লেখের বদলে $!A
  \lif !B$-এর !!{derivation}-এর শেষ পংক্তি এবং 2 নম্বর পংক্তির
  উল্লেখের বদলে $!B \lif !C$-এর !!{derivation}-এর শেষ পংক্তির
  উল্লেখ বসানো হয়।
\end{proof}

\begin{prob}
স্বতঃসিদ্ধগুলি থেকে !!{derivation}s প্রদর্শন করে দেখাও যে নিচেরগুলি সত্য:
\begin{enumerate}
\item $(!A \land !B) \lif (!B \land !A)$
\item $((!A \land !B) \lif !C) \lif (!A \lif (!B \lif !C))$
\item $\lnot(!A \lor !B) \lif \lnot !A$
\end{enumerate}
\end{prob}



\end{document}
